\documentclass[../../main.tex]{subfiles}% 注意这里的文件路径不能用 ./main.tex ,否则用latexmk编译子文件会报错
\graphicspath{{\subfix{./image/}}{\subfix{../../image/}}} % 指定图片引用路径,后续可以直接使用图片文件名
% 如果图片存放在上级目录,则使用 ../../image/ 作为备用路径

% 例如:
% \begin{figure}[H]
% \centering
% \includegraphics[width=0.3\textwidth]{图.png}
% \caption{}
% \label{figure:图}
% \end{figure}
% 注意:上述\label{}一定要放在\caption{}之后,否则引用图片序号会只会显示??.

\begin{document}
% SOURCE p=110; printed=95; object=ch03-euclidean-differentiation-title; region=upper
\section{欧氏空间上的微分}\label{section:ch03-euclidean-differentiation}
% SOURCE p=110; printed=95; object=ch03-euclidean-differentiation-intro; region=upper
Radon–Nikodym 定理给出了符号测度或复测度 $\nu$ 关于测度 $\mu$ 的“导数”的抽象概念。本节更深入地分析 $(X,\mathscr{M})=(\mathbb{R}^n,\mathscr{B}_{\mathbb{R}^n})$ 且 $\mu=m$ 为 Lebesgue 测度这一特殊情形。在这里，可以按下面的方式定义 $\nu$ 关于 $m$ 的逐点导数。
% SOURCE p=110; printed=95; object=ch03-pointwise-measure-derivative; region=upper-middle
\begin{definition}[测度的逐点导数]\label{definition:ch03-pointwise-measure-derivative}
以 $B(r,x)$ 表示 $\mathbb{R}^n$ 中以 $x$ 为中心、半径为 $r$ 的开球，考虑极限
\begin{align*}
F(x)=\lim_{r\to0}\frac{\nu(B(r,x))}{m(B(r,x))},
\end{align*}
只要该极限存在。
\end{definition}
\begin{remark}
也可以用在某种适当意义下以正规方式收缩到 $x$ 的其他集合代替球 $B(r,x)$；稍后将讨论这一点。
\end{remark}
% SOURCE p=110; printed=95; object=ch03-differentiation-motivation; region=middle
若 $\nu\ll m$，从而 $\mathrm{d}\nu=f\,\mathrm{d}m$，则 $\nu(B(r,x))/m(B(r,x))$ 就是 $f$ 在 $B(r,x)$ 上的平均值，因此可以期望 $F=f$，$m$-几乎处处成立。事实确实如此，只要对所有 $r,x$，$\nu(B(r,x))$ 都有限。从函数 $f$ 的角度看，这可视为微积分基本定理的推广：$f$ 的不定积分（即 $\nu$）的导数就是 $f$。

本节余下部分中，除非另有说明，“可积”“几乎处处”等术语都相对于 Lebesgue 测度。我们从一个本身也有独立意义的技术性引理开始。
% SOURCE p=110; printed=95; object=ch03-ball-covering; number=3.15; region=middle-lower
\begin{lemma}\label{lemma:ch03-ball-covering}
设 $\mathscr{C}$ 为 $\mathbb{R}^n$ 中的一族开球，$U=\bigcup_{B\in\mathscr{C}}B$。若 $c<m(U)$，则存在两两不交的 $B_1,\ldots,B_k\in\mathscr{C}$，使 $\sum_1^k m(B_j)>3^{-n}c$。
\end{lemma}
\begin{proof}
若 $c<m(U)$，由\cref{theorem:ch02-nd-regularity}，存在紧集 $K\subseteq U$，使 $m(K)>c$，而 $\mathscr{C}$ 中有限个球（记为 $A_1,\ldots,A_m$）覆盖 $K$。取 $B_1$ 为这些 $A_j$ 中最大的一个（即选取半径最大的球）。再取 $B_2$ 为与 $B_1$ 不交的 $A_j$ 中最大的一个，$B_3$ 为与 $B_1,B_2$ 都不交的 $A_j$ 中最大的一个，如此继续，直到这些 $A_j$ 都处理完。按此构造，若 $A_i$ 不在所选的 $B_j$ 中，则存在 $j$，使 $A_i\cap B_j\ne\varnothing$；若 $j$ 为满足此性质的最小整数，则 $A_i$ 的半径不超过 $B_j$ 的半径。所以 $A_i\subseteq B_j^*$，其中 $B_j^*$ 与 $B_j$ 同心，半径为 $B_j$ 的三倍。于是 $K\subseteq\bigcup_1^kB_j^*$，从而
\begin{align*}
c<m(K)\le\sum_1^k m(B_j^*)=3^n\sum_1^k m(B_j).
\end{align*}
\end{proof}
% SOURCE p=110,111; printed=95,96; object=ch03-locally-integrable; region=bottom/top
\begin{definition}[局部可积]\label{definition:ch03-locally-integrable}
若可测函数 $f:\mathbb{R}^n\to\mathbb{C}$ 满足：对每个有界可测集 $K\subseteq\mathbb{R}^n$，都有 $\int_K|f(x)|\,\mathrm{d}x<\infty$，则称 $f$（关于 Lebesgue 测度）\textbf{局部可积}。局部可积函数空间记作 $L^1_{\mathrm{loc}}$。
\end{definition}
% SOURCE p=111; printed=96; object=ch03-ball-average; region=top
\begin{definition}[球上的平均值]\label{definition:ch03-ball-average}
若 $f\in L^1_{\mathrm{loc}}$、$x\in\mathbb{R}^n$ 且 $r>0$，定义 $A_rf(x)$ 为 $f$ 在 $B(r,x)$ 上的平均值：
\begin{align*}
A_rf(x)\triangleq\frac{1}{m(B(r,x))}\int_{B(r,x)}f(y)\,\mathrm{d}y.
\end{align*}
\end{definition}
% SOURCE p=111; printed=96; object=ch03-ball-average-continuity; number=3.16; region=upper
\begin{lemma}\label{lemma:ch03-ball-average-continuity}
若 $f\in L^1_{\mathrm{loc}}$，则 $A_rf(x)$ 关于 $r$ 和 $x$ 联合连续（$r>0$，$x\in\mathbb{R}^n$）。
\end{lemma}
\begin{proof}
由\cref{section:ch02-polar-coordinates}中的结果，$m(B(r,x))=cr^n$，其中 $c=m(B(1,0))$，并且 $m(S(r,x))=0$，这里 $S(r,x)=\{y:|y-x|=r\}$。此外，当 $r\to r_0$、$x\to x_0$ 时，$\chi_{B(r,x)}\to\chi_{B(r_0,x_0)}$ 在 $\mathbb{R}^n\setminus S(r_0,x_0)$ 上逐点成立。所以 $\chi_{B(r,x)}\to\chi_{B(r_0,x_0)}$ 几乎处处成立；当 $r<r_0+\frac12$ 且 $|x-x_0|<\frac12$ 时，$|\chi_{B(r,x)}|\le\chi_{B(r_0+1,x_0)}$。由控制收敛定理，$\int_{B(r,x)}f(y)\,\mathrm{d}y$ 关于 $r,x$ 连续，因而 $A_rf(x)=c^{-1}r^{-n}\int_{B(r,x)}f(y)\,\mathrm{d}y$ 也如此。
\end{proof}
% SOURCE p=111; printed=96; object=ch03-hardy-littlewood-maximal; region=middle
\begin{definition}[Hardy–Littlewood 极大函数]\label{definition:ch03-hardy-littlewood-maximal}
若 $f\in L^1_{\mathrm{loc}}$，定义它的\textbf{Hardy–Littlewood 极大函数} $Hf$ 为
\begin{align*}
Hf(x)\triangleq\sup_{r>0}A_r|f|(x)
=\sup_{r>0}\frac{1}{m(B(r,x))}\int_{B(r,x)}|f(y)|\,\mathrm{d}y.
\end{align*}
\end{definition}
% SOURCE p=111; printed=96; object=ch03-maximal-measurable; region=middle
\begin{proposition}\label{proposition:ch03-maximal-measurable}
$Hf$ 为可测函数。
\end{proposition}
\begin{proof}
由\cref{lemma:ch03-ball-average-continuity}，对任意 $a\in\mathbb{R}$，
\begin{align*}
(Hf)^{-1}((a,\infty))=\bigcup_{r>0}(A_r|f|)^{-1}((a,\infty))
\end{align*}
为开集。
\end{proof}
% SOURCE p=111; printed=96; object=ch03-maximal-theorem; number=3.17; region=middle
\begin{theorem}[极大定理]\label{theorem:ch03-maximal}
存在常数 $C>0$，使对所有 $f\in L^1$ 和所有 $\alpha>0$，
\begin{align*}
m(\{x:Hf(x)>\alpha\})\le\frac{C}{\alpha}\int|f(x)|\,\mathrm{d}x.
\end{align*}
\end{theorem}
\begin{proof}
令 $E_\alpha=\{x:Hf(x)>\alpha\}$。对每个 $x\in E_\alpha$，可取 $r_x>0$，使 $A_{r_x}|f|(x)>\alpha$。球 $B(r_x,x)$ 覆盖 $E_\alpha$，所以由\cref{lemma:ch03-ball-covering}，若 $c<m(E_\alpha)$，则存在 $x_1,\ldots,x_k\in E_\alpha$，使球 $B_j=B(r_{x_j},x_j)$ 两两不交，且 $\sum_1^k m(B_j)>3^{-n}c$。于是
\begin{align*}
c<3^n\sum_1^k m(B_j)
\le\frac{3^n}{\alpha}\sum_1^k\int_{B_j}|f(y)|\,\mathrm{d}y
\le\frac{3^n}{\alpha}\int_{\mathbb{R}^n}|f(y)|\,\mathrm{d}y.
\end{align*}
令 $c\to m(E_\alpha)$，即得所需结论。
\end{proof}
% SOURCE p=111; printed=96; object=ch03-limsup-function; region=lower
有了这一工具，现在给出基本微分定理的三个依次加强的版本。证明中将用到实变量实值函数的上极限概念：
\begin{definition}[函数的上极限]\label{definition:ch03-function-limsup}
\begin{align*}
\limsup_{r\to R}\phi(r)
\triangleq\lim_{\varepsilon\to0}\sup_{0<|r-R|<\varepsilon}\phi(r)
=\inf_{\varepsilon>0}\sup_{0<|r-R|<\varepsilon}\phi(r).
\end{align*}
\end{definition}
还将用到下面这个容易验证的事实。
\begin{proposition}\label{proposition:ch03-limit-limsup-criterion}
\begin{align*}
\lim_{r\to R}\phi(r)=c\quad\Longleftrightarrow\quad\limsup_{r\to R}|\phi(r)-c|=0.
\end{align*}
\end{proposition}
% SOURCE p=112; printed=97; object=ch03-ball-differentiation; number=3.18; region=top
\begin{theorem}\label{theorem:ch03-ball-differentiation}
若 $f\in L^1_{\mathrm{loc}}$，则对几乎处处的 $x\in\mathbb{R}^n$，$\lim_{r\to0}A_rf(x)=f(x)$。
\end{theorem}
\begin{proof}
只需证明：对 $N\in\mathbb{N}$，$A_rf(x)\to f(x)$ 对几乎处处满足 $|x|\le N$ 的 $x$ 成立。但当 $|x|\le N$ 且 $r\le1$ 时，$A_rf(x)$ 只依赖于满足 $|y|\le N+1$ 的那些 $f(y)$，所以用 $f\chi_{B(N+1,0)}$ 代替 $f$，可以假设 $f\in L^1$。

给定 $\varepsilon>0$，由\cref{theorem:ch02-nd-lone-density}，可找到连续可积函数 $g$，使 $\int|g(y)-f(y)|\,\mathrm{d}y<\varepsilon$。由 $g$ 的连续性，对每个 $x\in\mathbb{R}^n$ 和 $\delta>0$，都存在 $r>0$，使当 $|y-x|<r$ 时，$|g(y)-g(x)|<\delta$，因而
\begin{align*}
|A_rg(x)-g(x)|=\frac{1}{m(B(r,x))}\left|\int_{B(r,x)}[g(y)-g(x)]\,\mathrm{d}y\right|<\delta.
\end{align*}
因此，当 $r\to0$ 时，对每个 $x$ 都有 $A_rg(x)\to g(x)$，所以
\begin{align*}
\limsup_{r\to0}|A_rf(x)-f(x)|
&=\limsup_{r\to0}|A_r(f-g)(x)+(A_rg-g)(x)+(g-f)(x)|\\
&\le H(f-g)(x)+0+|f-g|(x).
\end{align*}
于是，令
\begin{align*}
E_\alpha&=\left\{x:\limsup_{r\to0}|A_rf(x)-f(x)|>\alpha\right\},\\
F_\alpha&=\{x:|f-g|(x)>\alpha\},
\end{align*}
便有 $E_\alpha\subseteq F_{\alpha/2}\cup\{x:H(f-g)(x)>\alpha/2\}$。但 $(\alpha/2)m(F_{\alpha/2})\le\int_{F_{\alpha/2}}|f(x)-g(x)|\,\mathrm{d}x<\varepsilon$，所以由极大定理，
\begin{align*}
m(E_\alpha)\le\frac{2\varepsilon}{\alpha}+\frac{2C\varepsilon}{\alpha}.
\end{align*}
由于 $\varepsilon$ 任意，对所有 $\alpha>0$ 都有 $m(E_\alpha)=0$。而对所有 $x\notin\bigcup_1^\infty E_{1/n}$，均有 $\lim_{r\to0}A_rf(x)=f(x)$，证明完成。
\end{proof}
% SOURCE p=112; printed=97; object=ch03-ball-differentiation-equivalent; number=3.19; region=lower
这一结果可改写如下：若 $f\in L^1_{\mathrm{loc}}$，则
\begin{align}\label{align:ch03-ball-differentiation-equivalent}
\lim_{r\to0}\frac{1}{m(B(r,x))}\int_{B(r,x)}[f(y)-f(x)]\,\mathrm{d}y=0
\qquad\text{对几乎处处的 }x.
\end{align}
事实上，更强的结论也成立：即使把被积函数换成它的绝对值，\eqref{align:ch03-ball-differentiation-equivalent}仍成立。为此定义：
% SOURCE p=112; printed=97; object=ch03-lebesgue-set; region=bottom
\begin{definition}[Lebesgue 集]\label{definition:ch03-lebesgue-set}
$f$ 的\textbf{Lebesgue 集} $L_f$ 定义为
\begin{align*}
L_f\triangleq\left\{x:\lim_{r\to0}\frac{1}{m(B(r,x))}\int_{B(r,x)}|f(y)-f(x)|\,\mathrm{d}y=0\right\}.
\end{align*}
\end{definition}
% SOURCE p=113; printed=98; object=ch03-lebesgue-set-full; number=3.20; region=top
\begin{theorem}\label{theorem:ch03-lebesgue-set-full}
若 $f\in L^1_{\mathrm{loc}}$，则 $m((L_f)^c)=0$。
\end{theorem}
\begin{proof}
对每个 $c\in\mathbb{C}$，将\cref{theorem:ch03-ball-differentiation}用于 $g_c(x)=|f(x)-c|$，得到：在一个 Lebesgue 零集 $E_c$ 之外，有
\begin{align*}
\lim_{r\to0}\frac{1}{m(B(r,x))}\int_{B(r,x)}|f(y)-c|\,\mathrm{d}y=|f(x)-c|.
\end{align*}
取 $\mathbb{C}$ 的一个可数稠密子集 $D$，令 $E=\bigcup_{c\in D}E_c$。则 $m(E)=0$。若 $x\notin E$，对任意 $\varepsilon>0$，可取 $c\in D$，使 $|f(x)-c|<\varepsilon$，于是 $|f(y)-f(x)|<|f(y)-c|+\varepsilon$，从而
\begin{align*}
\limsup_{r\to0}\frac{1}{m(B(r,x))}\int_{B(r,x)}|f(y)-f(x)|\,\mathrm{d}y
\le|f(x)-c|+\varepsilon<2\varepsilon.
\end{align*}
由 $\varepsilon$ 的任意性，结论得证。
\end{proof}
% SOURCE p=113; printed=98; object=ch03-shrink-nicely; region=middle
最后，考虑比球更一般的集合族。
\begin{definition}[正规收缩]\label{definition:ch03-shrink-nicely}
$\mathbb{R}^n$ 的 Borel 子集族 $\{E_r\}_{r>0}$ 称为\textbf{正规收缩}到 $x\in\mathbb{R}^n$，若：
\begin{enumerate}
\item 对每个 $r$，$E_r\subseteq B(r,x)$。
\item 存在与 $r$ 无关的常数 $\alpha>0$，使 $m(E_r)>\alpha m(B(r,x))$。
\end{enumerate}
\end{definition}
\begin{remark}
集合 $E_r$ 不必包含 $x$ 本身。
\end{remark}
% SOURCE p=113; printed=98; object=ch03-shrink-nicely-example; region=middle
\begin{example}\label{example:ch03-shrink-nicely}
若 $U$ 为 $B(1,0)$ 的任意 Borel 子集，$m(U)>0$，且 $E_r=\{x+ry:y\in U\}$，则 $\{E_r\}$ 正规收缩到 $x$。
\end{example}
下面给出微分定理的最终形式。
% SOURCE p=113; printed=98; object=ch03-lebesgue-differentiation; number=3.21; region=lower
\begin{theorem}[Lebesgue 微分定理]\label{theorem:ch03-lebesgue-differentiation}
设 $f\in L^1_{\mathrm{loc}}$。对 $f$ 的 Lebesgue 集中的每个 $x$（特别地，对几乎处处的 $x$），以及每个正规收缩到 $x$ 的集合族 $\{E_r\}_{r>0}$，都有
\begin{align*}
\lim_{r\to0}\frac{1}{m(E_r)}\int_{E_r}|f(y)-f(x)|\,\mathrm{d}y&=0,\\
\lim_{r\to0}\frac{1}{m(E_r)}\int_{E_r}f(y)\,\mathrm{d}y&=f(x).
\end{align*}
\end{theorem}
\begin{proof}
对某个 $\alpha>0$，有
\begin{align*}
\frac{1}{m(E_r)}\int_{E_r}|f(y)-f(x)|\,\mathrm{d}y
&\le\frac{1}{m(E_r)}\int_{B(r,x)}|f(y)-f(x)|\,\mathrm{d}y\\
&\le\frac{1}{\alpha m(B(r,x))}\int_{B(r,x)}|f(y)-f(x)|\,\mathrm{d}y.
\end{align*}
所以第一个等式由\cref{theorem:ch03-lebesgue-set-full}得出；将第二个等式写成\eqref{align:ch03-ball-differentiation-equivalent}的形式，立即看出第一个等式蕴含第二个。
\end{proof}
% SOURCE p=114; printed=99; object=ch03-regular-borel-measure; region=top
现在回到测度的研究。
\begin{definition}[正则 Borel 测度]\label{definition:ch03-regular-borel-measure}
$\mathbb{R}^n$ 上的 Borel 测度 $\nu$ 称为\textbf{正则的}，若：
\begin{enumerate}[label=(\roman*)]
\item\label{item:ch03-regular-finite-compact} 对每个紧集 $K$，$\nu(K)<\infty$。
\item\label{item:ch03-regular-outer} 对每个 $E\in\mathscr{B}_{\mathbb{R}^n}$，$\nu(E)=\inf\{\nu(U):U\text{ 为开集},E\subseteq U\}$。
\end{enumerate}
\end{definition}
\begin{remark}
条件\subcref{definition:ch03-regular-borel-measure}{item:ch03-regular-outer}实际上由条件\subcref{definition:ch03-regular-borel-measure}{item:ch03-regular-finite-compact}蕴含。对 $n=1$，这由\cref{theorem:ch01-stieltjes-measure,theorem:ch01-measure-properties}得出；对于任意 $n$，将在\cref{section:ch07-regular-borel-measures}中证明。暂且将条件\subcref{definition:ch03-regular-borel-measure}{item:ch03-regular-outer}明确作为假设。
\end{remark}
% SOURCE p=114; printed=99; object=ch03-regular-sigma-finite; region=upper
\begin{proposition}\label{proposition:ch03-regular-sigma-finite}
每个正则测度都是 $\sigma$-有限的。
\end{proposition}
\begin{proof}
由\subcref{definition:ch03-regular-borel-measure}{item:ch03-regular-finite-compact}即可看出。
\end{proof}
% SOURCE p=114; printed=99; object=ch03-regular-signed-complex; region=upper
\begin{definition}[正则符号测度与复测度]\label{definition:ch03-regular-signed-complex}
若 $|\nu|$ 正则，则称符号或复 Borel 测度 $\nu$ 为\textbf{正则的}。
\end{definition}
% SOURCE p=114; printed=99; object=ch03-regular-density-example; region=upper-middle
\begin{example}\label{example:ch03-regular-density}
若 $f\in L^+(\mathbb{R}^n)$，则测度 $f\,\mathrm{d}m$ 正则，当且仅当 $f\in L^1_{\mathrm{loc}}$。
\end{example}
\begin{proof}
条件 $f\in L^1_{\mathrm{loc}}$ 显然等价于\subcref{definition:ch03-regular-borel-measure}{item:ch03-regular-finite-compact}。若它成立，可以直接验证\subcref{definition:ch03-regular-borel-measure}{item:ch03-regular-outer}如下。设 $E$ 为有界 Borel 集。给定 $\delta>0$，由\cref{theorem:ch02-nd-regularity}，存在有界开集 $U\supseteq E$，使 $m(U)<m(E)+\delta$，从而 $m(U\setminus E)<\delta$。于是，给定 $\varepsilon>0$，由\cref{corollary:ch03-integral-absolute-continuity}，存在开集 $U\supseteq E$，使 $\int_{U\setminus E}f\,\mathrm{d}m<\varepsilon$，从而 $\int_U f\,\mathrm{d}m<\int_E f\,\mathrm{d}m+\varepsilon$。对于无界的 $E$，将其写成 $E=\bigcup_1^\infty E_j$，其中每个 $E_j$ 有界，再取开集 $U_j\supseteq E_j$，使 $\int_{U_j\setminus E_j}f\,\mathrm{d}m<\varepsilon2^{-j}$，即可得到结论。
\end{proof}
% SOURCE p=114; printed=99; object=ch03-measure-differentiation; number=3.22; region=middle
\begin{theorem}\label{theorem:ch03-measure-differentiation}
设 $\nu$ 为 $\mathbb{R}^n$ 上的正则符号或复 Borel 测度，$\mathrm{d}\nu=\mathrm{d}\lambda+f\,\mathrm{d}m$ 为其 Lebesgue–Radon–Nikodym 表示。则对 $m$-几乎处处的 $x\in\mathbb{R}^n$，以及每个正规收缩到 $x$ 的集合族 $\{E_r\}_{r>0}$，都有
\begin{align*}
\lim_{r\to0}\frac{\nu(E_r)}{m(E_r)}=f(x).
\end{align*}
\end{theorem}
\begin{proof}
容易验证 $\mathrm{d}|\nu|=\mathrm{d}|\lambda|+|f|\,\mathrm{d}m$，所以 $\nu$ 的正则性蕴含 $\lambda$ 与 $f\,\mathrm{d}m$ 的正则性（\cref{exercise:ch03-regular-singular-components}）。特别地，$f\in L^1_{\mathrm{loc}}$。由\cref{theorem:ch03-lebesgue-differentiation}，只需证明：若 $\lambda$ 正则且 $\lambda\perp m$，则对 $m$-几乎处处的 $x$，当 $E_r$ 正规收缩到 $x$ 时，$\lambda(E_r)/m(E_r)\to0$。还只需取 $E_r=B(r,x)$，并假设 $\lambda$ 为正测度，因为对某个 $\alpha>0$，有
\begin{align*}
\left|\frac{\lambda(E_r)}{m(E_r)}\right|
\le\frac{|\lambda|(E_r)}{m(E_r)}
\le\frac{|\lambda|(B(r,x))}{m(E_r)}
\le\frac{|\lambda|(B(r,x))}{\alpha m(B(r,x))}.
\end{align*}
于是，假设 $\lambda\ge0$，取 Borel 集 $A$，使 $\lambda(A)=m(A^c)=0$，令
\begin{align*}
F_k=\left\{x\in A:\limsup_{r\to0}\frac{\lambda(B(r,x))}{m(B(r,x))}>\frac1k\right\}.
\end{align*}
我们将证明对所有 $k$，$m(F_k)=0$，从而完成证明。

论证与极大定理的证明相似。由 $\lambda$ 的正则性，给定 $\varepsilon>0$，存在开集 $U_\varepsilon\supseteq A$，使 $\lambda(U_\varepsilon)<\varepsilon$。每个 $x\in F_k$ 都是某个球 $B_x\subseteq U_\varepsilon$ 的中心，且 $\lambda(B_x)>k^{-1}m(B_x)$。由\cref{lemma:ch03-ball-covering}，
% SOURCE p=115; printed=100; object=ch03-measure-differentiation-proof-cont; region=top
若 $V_\varepsilon=\bigcup_{x\in F_k}B_x$ 且 $c<m(V_\varepsilon)$，则存在 $x_1,\ldots,x_J$，使 $B_{x_1},\ldots,B_{x_J}$ 两两不交，且
\begin{align*}
c<3^n\sum_1^Jm(B_{x_j})
&\le3^nk\sum_1^J\lambda(B_{x_j})\le3^nk\lambda(V_\varepsilon)\\
&\le3^nk\lambda(U_\varepsilon)\le3^nk\varepsilon.
\end{align*}
故 $m(V_\varepsilon)\le3^nk\varepsilon$；由 $F_k\subseteq V_\varepsilon$ 和 $\varepsilon$ 的任意性，$m(F_k)=0$。
\end{proof}
% SOURCE p=115; printed=100; object=ch03-euclidean-exercises-title; region=upper
\subsection*{习题}
% SOURCE p=115; printed=100; object=ch03-ex22; number=22; region=upper
\begin{exercise}\label{exercise:ch03-maximal-sharpness}
若 $f\in L^1(\mathbb{R}^n)$ 且 $f\ne0$，则存在 $C,R>0$，使当 $|x|>R$ 时，$Hf(x)\ge C|x|^{-n}$。因此，当 $\alpha$ 很小时，$m(\{x:Hf(x)>\alpha\})\ge C'/\alpha$，所以极大定理中的估计本质上是尖锐的。
\end{exercise}
% SOURCE p=115; printed=100; object=ch03-ex23; number=23; region=upper-middle
\begin{exercise}\label{exercise:ch03-uncentered-maximal}
\begin{definition*}[极大函数的一个变体]
Hardy–Littlewood 极大函数的一个有用变体是
\begin{align*}
H^*f(x)\triangleq\sup\left\{\frac1{m(B)}\int_B|f(y)|\,\mathrm{d}y:B\text{ 为球且 }x\in B\right\}.
\end{align*}
\end{definition*}
证明 $Hf\le H^*f\le2^nHf$。
\end{exercise}
% SOURCE p=115; printed=100; object=ch03-ex24; number=24; region=middle
\begin{exercise}\label{exercise:ch03-continuity-lebesgue-point}
若 $f\in L^1_{\mathrm{loc}}$，且 $f$ 在 $x$ 处连续，则 $x$ 属于 $f$ 的 Lebesgue 集。
\end{exercise}
% SOURCE p=115; printed=100; object=ch03-ex25; number=25; region=middle
\begin{exercise}\label{exercise:ch03-set-density}
\begin{definition*}[集合在一点的密度]
若 $E$ 为 $\mathbb{R}^n$ 中的 Borel 集，则 $E$ 在 $x$ 处的\textbf{密度} $D_E(x)$ 定义为
\begin{align*}
D_E(x)\triangleq\lim_{r\to0}\frac{m(E\cap B(r,x))}{m(B(r,x))},
\end{align*}
只要该极限存在。
\end{definition*}
\begin{enumerate}
\item 证明：对几乎处处的 $x\in E$，$D_E(x)=1$；对几乎处处的 $x\in E^c$，$D_E(x)=0$。
\item 给出 $E,x$ 的例子，使 $D_E(x)$ 等于给定的 $\alpha\in(0,1)$，或使 $D_E(x)$ 不存在。
\end{enumerate}
\end{exercise}
% SOURCE p=115; printed=100; object=ch03-ex26; number=26; region=middle-lower
\begin{exercise}\label{exercise:ch03-regular-singular-components}
若 $\lambda,\mu$ 为 $\mathbb{R}^n$ 上相互奇异的正 Borel 测度，且 $\lambda+\mu$ 正则，则 $\lambda,\mu$ 也正则。
\end{exercise}
\end{document}
